\documentclass[11pt,letterpaper]{article}
\usepackage[utf8]{inputenc}
\usepackage[T1]{fontenc}
\usepackage[spanish,es-noshorthands,es-tabla]{babel}
\usepackage[margin=2.3cm]{geometry}
\usepackage{lmodern}
\usepackage{booktabs}
\usepackage{longtable}
\usepackage{array}
\usepackage{xcolor}
\usepackage[hidelinks]{hyperref}
\usepackage{enumitem}
\setlist{nosep,leftmargin=*}
\newcolumntype{L}[1]{>{\raggedright\arraybackslash}p{#1}}
\newcommand{\ruta}[1]{\texttt{\small #1}}
\sloppy
\renewcommand{\arraystretch}{1.15}

\title{Tukit: informe de control de calidad\\[4pt]
\large Informe de pruebas según la estructura de IEEE 29119-3\\
Proyecto de grado ISIS2610, grupo 19, Pontificia Universidad Javeriana}
\author{Equipo Tukit}
\date{1 de noviembre de 2026 (re-ejecución del 2 de noviembre)}

\begin{document}
\maketitle

\section{Propósito y alcance}

Este informe reúne la evidencia de calidad de Tukit en la rama \ruta{validacion/\allowbreak{}oe5}, commit \ruta{a664810}
(1 de noviembre de 2026, 23:33, hora local UTC$-$8), que es la versión congelada para la ronda sumativa de
usabilidad del OE5. Sigue la estructura del informe de finalización de pruebas de IEEE 29119-3: contexto,
pruebas ejecutadas por nivel, desviaciones, hallazgos, trazabilidad hacia los requisitos y lo que queda pendiente.

La regla que gobierna el documento es simple: solo se reporta una prueba que de verdad se ejecutó, con su
fecha, su comando y su resultado. Lo que no se corrió aparece como ``no ejecutado'', con la razón. Las cifras
de producto (servicios, tablas, pesos, plazos legales) se toman de la hoja de hechos canónicos
(\ruta{HECHOS\_\allowbreak{}CANONICOS\_\allowbreak{}DOCS.md}); las de pruebas salen de los registros que acompañan este informe, en la
carpeta \ruta{QC/\allowbreak{}evidencia/}.

Para leer las tablas conviene distinguir tres tipos de evidencia. Una \emph{prueba automatizada} es código que
se ejecuta y falla si el comportamiento cambia. Una \emph{prueba de extremo a extremo} (E2E) recorre el sistema
vivo por HTTP, como lo haría un acudiente. Una \emph{inspección} es la lectura dirigida del código o de una
página contra una regla escrita, sin un caso automatizado detrás.

\section{Entorno de prueba}

\begin{itemize}
  \item Máquina de desarrollo macOS (Apple Silicon), sin Docker. PostgreSQL 16 con pgvector 0.8.1,
        Redis y RabbitMQ instalados con Homebrew y atados a \ruta{127.0.0.1}.
  \item Python 3.13.11 en el entorno \ruta{.venv} del monorepo; Node.js 22.23.4.
  \item Servicios vivos: puerta de enlace en el puerto 8010 y los cinco servicios de dominio en 8011 a 8015;
        la web del acudiente en el puerto 3000.
  \item Pasarela Wompi en modo simulado (\ruta{WOMPI\_\allowbreak{}MODO} por omisión); no hay llaves de sandbox.
\end{itemize}

Todas las ejecuciones de este informe ocurrieron el 1 de noviembre de 2026 entre las 23:35 y las 23:45 hora
local (2 de noviembre, 07:35 a 07:45 UTC), salvo la E2E de la coordinación, que se corrió el mismo día por un
túnel público hacia la web local, y los artefactos de evaluación de la sección~\ref{sec:artefactos}, que se
citan con su propia fecha.

\section{Resumen por nivel}

La tabla~\ref{tab:resumen} condensa los resultados. Ningún caso automatizado falló. Los problemas encontrados
están en otro lado: en el análisis estático, en la prueba de humo, en las dependencias y en una ruta que no
comprueba quién es el dueño del recurso (sección~\ref{sec:hallazgos}).

\begin{table}[h]
\centering\small
\caption{Resultados por nivel de prueba (1 de noviembre de 2026).}\label{tab:resumen}
\begin{tabular}{L{4.6cm}L{4.4cm}rrrL{2.2cm}}
\toprule
Nivel & Comando & Pasan & Fallan & Saltadas & Estado \\
\midrule
Fronteras de arquitectura (Python) & \ruta{make fronteras} & 75 & 0 & 1 & pasa \\
Unitarias e integración (Python) & \ruta{make prueba} & 2.481 & 0 & 1 & pasa \\
Contra RabbitMQ real & \ruta{make prueba-colas} & 6 & 0 & 0 & pasa \\
Contra Redis real & \ruta{make prueba-redis} & 11 & 0 & 0 & pasa \\
Web del acudiente (TS) & \ruta{npm test} en \ruta{apps/\allowbreak{}web} & 82 & 0 & 0 & pasa \\
Consola admin (TS) & \ruta{node --test} en \ruta{apps/\allowbreak{}admin} & 57 & 0 & 0 & pasa \\
Contrato compartido (TS) & \ruta{npm test} en \ruta{libs/\allowbreak{}ts/\allowbreak{}tukit-contrato} & 67 & 0 & 0 & pasa \\
Tipos TS & \ruta{tsc --noEmit} por paquete & 2 de 3 & 1 & & falla (1 error) \\
Estilo Python & \ruta{make lint} & & 62 & & falla \\
Estilo TS & \ruta{npm run lint} & admin & web (1) & & falla \\
Humo sobre servicios vivos & \ruta{make humo} & & 1 & & falla (ver H2) \\
Humo corregido (20 comprobaciones) & \ruta{humo\_\allowbreak{}traza.py} & 20 & 0 & 0 & pasa \\
E2E del acudiente & scripts 
uta{e2e*.sh} & 12 pasos & 0 & & pasa \\
Acceso cruzado (IDOR) & \ruta{idor.sh} & 5 & 3 & & hallazgo H4 \\
Cobertura de código & no hay \ruta{pytest-cov} & & & & no ejecutado \\
\bottomrule
\end{tabular}
\end{table}

En total, 2.573 pruebas automatizadas de Python pasan, 2 se saltan y ninguna falla; las 206 de TypeScript
pasan todas.

\section{Pruebas automatizadas}

\subsection{Python}

El monorepo corre la suite desde la raíz con \ruta{pytest} y una lista de \ruta{pythonpath} declarada en
\ruta{pyproject.toml}, de modo que cada servicio se prueba como paquete propio sin instalarlo. El objetivo
\ruta{make prueba} corre primero \ruta{make fronteras} (las funciones de aptitud de arquitectura, que fallan
si un servicio importa el código o la base de otro) y luego el resto de la suite sin los marcadores
\ruta{rabbitmq} y \ruta{redis}. Esos 17 casos se corrieron aparte, contra el broker y el Redis locales.

\begin{table}[h]
\centering\small
\caption{Pruebas de Python por paquete (\ruta{pytest <paquete> -m "not rabbitmq and not redis"}).}
\begin{tabular}{lrrl}
\toprule
Paquete & Pasan & Saltadas & Observación \\
\midrule
\ruta{arquitectura} (fronteras) & 75 & 1 & salta \ruta{ingesta}: aún sin código \\
\ruta{libs} (\ruta{tukit\_\allowbreak{}comun}) & 169 & 0 & 14 casos van a colas y Redis \\
\ruta{servicios/\allowbreak{}identidad} & 220 & 0 & \\
\ruta{servicios/\allowbreak{}catalogo} & 813 & 0 & \\
\ruta{servicios/\allowbreak{}inteligencia} & 535 & 0 & \\
\ruta{servicios/\allowbreak{}comercio} & 271 & 0 & \\
\ruta{servicios/\allowbreak{}operacion} & 110 & 0 & \\
\ruta{servicios/\allowbreak{}gateway} & 222 & 0 & \\
\ruta{ensayos} & 141 & 1 & salta \ruta{nl2c\_\allowbreak{}vivo}: requiere el modelo \\
\midrule
Total & 2.556 & 2 & más 17 de colas y Redis: 2.573 \\
\bottomrule
\end{tabular}
\end{table}

Las dos pruebas saltadas lo están a propósito. La de \ruta{arquitectura/\allowbreak{}test\_\allowbreak{}fronteras.py:614} espera a que el
desplegable de ingesta tenga código, y la de \ruta{ensayos/\allowbreak{}nl2c/\allowbreak{}test\_\allowbreak{}nl2c\_\allowbreak{}ensayo.py:251} solo corre con
\ruta{-m nl2c\_\allowbreak{}vivo}, porque llama al modelo real; su equivalente ya ejecutado son las rondas del ensayo NL2C.

\textbf{Cobertura.} No ejecutada: el entorno no tiene \ruta{pytest-cov} y no hay configuración de cobertura en
\ruta{pyproject.toml}. No se instaló nada para este informe. Por eso el requisito MT-01 (cobertura de al menos
80~\% como condición de integración) queda sin verificar.

\subsection{TypeScript}

Los tres paquetes de TypeScript usan el corredor nativo de Node (\ruta{node --test "test/\allowbreak{}**/\allowbreak{}*.test.ts"}).
La consola admin no declara un script \ruta{test} en su \ruta{package.json}, así que se invocó el mismo
comando directamente. La web pasa 82 de 82, la consola admin 57 de 57 y el contrato compartido 67 de 67.

La verificación de tipos (\ruta{tsc --noEmit}) pasa en \ruta{apps/\allowbreak{}web} y en \ruta{apps/\allowbreak{}admin}, pero falla en
\ruta{libs/\allowbreak{}ts/\allowbreak{}tukit-contrato} con un error (hallazgo H1).

\subsection{Análisis estático}

\ruta{make lint} falla con 62 avisos de \ruta{ruff}: 47 líneas largas (E501), 9 redefiniciones (F811), 3
importaciones sin ordenar, 2 nombres ambiguos y 1 importación sin usar; además, 16 archivos no cumplen el
formato. Siete de las redefiniciones están en \ruta{servicios/\allowbreak{}comercio/\allowbreak{}tests/\allowbreak{}test\_\allowbreak{}admin\_\allowbreak{}pedido.py} y son el
patrón conocido de importar una \emph{fixture} y recibirla como argumento; las dos de
\ruta{servicios/\allowbreak{}operacion/\allowbreak{}tukit\_\allowbreak{}operacion/\allowbreak{}remoto.py} (líneas 45 y 92) son código de producción y merecen
revisión. En TypeScript, \ruta{eslint} pasa en la consola admin y marca un error en la web:
\ruta{apps/\allowbreak{}web/\allowbreak{}src/\allowbreak{}app/\allowbreak{}privacidad/\allowbreak{}AccionesPrivacidad.tsx:85} usa \ruta{<a href="/\allowbreak{}">} en vez de \ruta{<Link>}.
Ninguno de estos avisos cambia el comportamiento observado en las pruebas.

\section{Prueba de humo}

La prueba de humo (\ruta{arquitectura/\allowbreak{}humo.py}) golpea los endpoints vivos a través de la puerta de enlace y
falla si alguno no trae datos reales. Recorre el catálogo, los colegios y las listas, lanza un trabajo,
espera el pipeline, lo aprueba como operador, verifica que los kits se publiquen, revisa pedidos y casos, y
abre un caso PQRS. Escribe datos a propósito: deja un trabajo de prueba aprobado y un caso abierto por cada
corrida.

\begin{enumerate}
  \item \ruta{make humo} falló en el primer paso autenticado: \ruta{/\allowbreak{}v1/\allowbreak{}sesiones} respondió 401
        \ruta{credenciales\_\allowbreak{}invalidas}. El script tiene la clave de demostración escrita en el código, y esa
        clave se rotó el 28 de octubre para las sesiones del OE5.
  \item Con la clave rotada (leída del llavero de macOS en tiempo de ejecución, sin escribirla en ningún
        archivo) el humo avanzó 19 comprobaciones y falló en \ruta{GET /\allowbreak{}v1/\allowbreak{}operacion/\allowbreak{}casos} con 401
        \ruta{sesion\_\allowbreak{}invalida}. Desde el cambio del 1 de noviembre la bandeja de operación es solo para el
        equipo, y el script la pide sin token.
  \item Con la bandeja pedida con la sesión del operador, pasaron las 20 comprobaciones: 5 procesos arriba,
        146 conceptos, 1.078 productos con precio, 10 tiendas, 3 colegios, 19 listas, un trabajo que llegó a
        \ruta{espera\_\allowbreak{}revision} con sus cuatro etapas en \ruta{ok}, la cola admin cerrada sin token (401),
        la aprobación y la publicación de 3 kits (Ahorrador, A tu Medida y Premium, con A tu Medida
        recomendado), pedidos cerrados sin token (401) y el alta de un caso PQRS (201, vence el 17 de
        noviembre).
\end{enumerate}

La conclusión es que el sistema vivo responde como se espera y que el defecto está en el guion de humo, que
quedó desactualizado frente a dos cambios deliberados (hallazgo H2). Los conteos del catálogo son los de la
base de esta noche y cambian mientras avanza la curaduría; no deben citarse como cifras finales del catálogo.

\section{Prueba de extremo a extremo del acudiente}

La coordinación ejecutó el recorrido del acudiente el 1 de noviembre de 2026 con los guiones
\ruta{e2e.sh}, \ruta{e2e2.sh} y \ruta{e2e3.sh}, contra las rutas \ruta{/\allowbreak{}api/\allowbreak{}*} de la web y con la cookie real
de sesión. Se hizo por HTTP y no con un navegador automatizado porque Playwright cierra páginas de forma
impredecible en esta máquina. La tabla~\ref{tab:e2e} recoge lo observado.

\begin{longtable}{L{5.2cm}L{9.6cm}}
\caption{Pasos de la E2E del acudiente (1 de noviembre de 2026).}\label{tab:e2e}\\
\toprule
Paso & Resultado observado \\
\midrule\endhead
Registro con hijo y aceptación de términos & cuenta creada; la cookie de sesión queda puesta \\
Aplicar kit canónico, perfil Fan, personaje Spider-Man & publicado con 72 líneas, 10 productos del personaje;
2 líneas en \$0, ambas explicadas como duplicados \\
Carrusel de opciones de una línea & 42 opciones \\
Descartar una línea & \ruta{PUT} 200; el total baja de \$779.037 a \$757.038 \\
Asistente NL2C, proponer & devuelve una \ruta{PropuestaCambio} \\
Ordenar & pedido en \ruta{pendiente\_\allowbreak{}pago} \\
Pago simulado y sincronización & pedido en \ruta{pagado} \\
Caso PQRS & 201, vence en 48 horas; aparece en el listado del acudiente \\
Retracto sin pagar & 409 \\
Retracto después de pagar & 201; pedido \ruta{cancelado}, monto \$779.037 \\
Exportar mis datos & 200 \\
Páginas dashboard, privacidad, pedido, kit y legal & 200 \\
Pago real con el sandbox de Wompi & \textbf{no ejecutado}: las llaves del sandbox están pendientes \\
\bottomrule
\end{longtable}

\section{Seguridad}

\subsection{Acceso cruzado entre cuentas (IDOR)}

Una vulnerabilidad IDOR (\emph{insecure direct object reference}) aparece cuando basta conocer el
identificador de un recurso para leerlo, aunque sea de otra persona. Para probarlo se registró una cuenta
nueva de prueba (``atacante'') y, con su cookie, se pidieron por la web local los recursos de la cuenta de la
E2E. Todas las peticiones fueron de lectura; no se intentó escribir sobre recursos ajenos. El guion está en
\ruta{QC/\allowbreak{}evidencia/\allowbreak{}idor.sh} y la salida en \ruta{QC/\allowbreak{}evidencia/\allowbreak{}idor.log}.

\begin{table}[h]
\centering\small
\caption{Pruebas de acceso cruzado (1 de noviembre de 2026, 23:41).}
\begin{tabular}{L{6.3cm}cL{6.2cm}}
\toprule
Petición & Código & Lectura \\
\midrule
Pedido ajeno con la cookie del atacante & 404 & correcto: igual que un pedido inexistente \\
Pedido ajeno sin cookie & 401 & correcto \\
Pedido inventado & 404 & correcto \\
Casos del atacante & 200, 0 casos & correcto: no ve los casos de la otra cuenta \\
Exportación de datos del atacante & 200 & correcto: no contiene el pedido ni el trabajo ajenos \\
Trabajo ajeno con la cookie del atacante & 200 & \textbf{hallazgo H4}: devuelve el kit completo \\
Trabajo ajeno sin cookie & 200 & \textbf{hallazgo H4} \\
Opciones de una línea del trabajo ajeno & 200 & \textbf{hallazgo H4} \\
\bottomrule
\end{tabular}
\end{table}

No se probó por identificador el caso PQRS ajeno porque no se tenía su identificador; ese control lo cubre la
prueba automatizada \ruta{test\_\allowbreak{}un\_\allowbreak{}caso\_\allowbreak{}ajeno\_\allowbreak{}es\_\allowbreak{}403\_\allowbreak{}al\_\allowbreak{}leerlo\_\allowbreak{}y\_\allowbreak{}al\_\allowbreak{}comentarlo}, que pasa.

\subsection{Secretos en el repositorio}

Se corrió \ruta{gitleaks} 8.30.1 con \ruta{--redact}, de modo que ningún valor se imprimió ni quedó en los
resúmenes. Sobre el historial (570 commits) encontró 2 coincidencias de la regla \ruta{github-pat}, en dos
archivos de traspaso de sesión de \ruta{.opencode/\allowbreak{}handoff/\allowbreak{}history/} que entraron en el commit \ruta{d0c19b0f}
(4 de octubre de 2026). Esos archivos ya no están en el árbol de trabajo, pero siguen en el historial
(hallazgo H3).

Sobre el árbol de trabajo (\ruta{--no-git}) hay 1.482 coincidencias: 1.471 están en
\ruta{datos/\allowbreak{}catalogo/\allowbreak{}cache\_\allowbreak{}tiendas/}, que son respuestas guardadas de la tienda Éxito con tokens JWT
propios de esa tienda y que \ruta{.gitignore} excluye; 10 están en dependencias y artefactos de compilación; 1
está en \ruta{.env}, que también está excluido de git y no se leyó. Ninguna de estas 1.482 está versionada.

\subsection{Dependencias}

\ruta{npm audit --omit=dev} reporta lo mismo para la web y la consola admin, porque comparten el
\ruta{package-lock.json} del monorepo: 1 vulnerabilidad crítica y 2 altas. La crítica es \ruta{next} 16.3.4,
dentro del rango afectado 16.0.0 a 16.4.0, con dos avisos (ejecución remota de código en
\ruta{next/\allowbreak{}og ImageResponse} y una fuga del modo borrador con \ruta{use cache}); la corrección es 16.4.3. Las
altas son \ruta{sharp} (dependencia de \ruta{librsvg}, CVE-2026-96889) y \ruta{source-map-js} (denegación de
servicio). Una búsqueda en \ruta{apps/} no encontró usos de \ruta{next/\allowbreak{}og}, \ruta{ImageResponse} ni
\ruta{use cache}, lo que reduce la exposición pero no la elimina (hallazgo H5).

La auditoría de dependencias de Python \textbf{no se ejecutó}: \ruta{pip-audit} no está instalado y no se
instaló nada para este informe.

\section{Artefactos de evaluación ya ejecutados}\label{sec:artefactos}

Estas evaluaciones se corrieron antes de este informe y no se repitieron; se citan desde sus archivos.

\subsection{Ensayo NL2C}

El ensayo mide si el intérprete de lenguaje natural (\ruta{gemini-3.5-flash-lite}, temperatura 0) traduce
bien las frases del acudiente a la estructura $R$ y si resiste inyecciones de instrucciones. La tasa de
éxito de ataque (ASR, \emph{attack success rate}) cuenta los ataques que llegan al padre, no los que el modelo
intenta. Fuente: \ruta{docs/\allowbreak{}informes/\allowbreak{}nl2c/\allowbreak{}tablas\_\allowbreak{}ronda2.md} y \ruta{tablas\_\allowbreak{}ronda3.md}.

\begin{table}[h]
\centering\small
\caption{Rondas 2 y 3 del ensayo NL2C.}
\begin{tabular}{lll}
\toprule
Medida & Ronda 2 & Ronda 3 \\
\midrule
Corrida (UTC) y commit & 2026-11-02 02:36, \ruta{0373d5e} & 2026-11-02 03:12, \ruta{bf8b7d5} \\
Tarea completa & 57 de 70 (81,4~\%) & 58 de 70 (82,9~\%) \\
ASR & 0 de 81 & 0 de 81 \\
Latencia p50 / p95 & 999 / 1.344 ms & 1.000 / 1.177 ms \\
Tareas cumplidas bajo ataque & 36 de 60 & 36 de 60 \\
Controles sin ataque sobre el kit & 12 de 12 & 12 de 12 \\
\bottomrule
\end{tabular}
\end{table}

La ronda 3 usa un criterio modificado (cuenta como fallo un descarte sin confirmación del padre), así que las
dos rondas no son directamente comparables. La hoja de hechos canónicos todavía marca estas cifras como
pendientes de confirmar y pide no escribirlas como resultado final en los documentos de requisitos; aquí se
reportan como evidencia de prueba, con su fecha. Que el modelo cumpla la tarea legítima solo en 36 de 60
intentos cuando la frase trae un ataque es una limitación de utilidad, no de seguridad: ninguno de esos ataques
llegó al padre.

\subsection{Calibración del optimizador}

Fuente: \ruta{docs/\allowbreak{}informes/\allowbreak{}calibracion/} (README y JSON). La configuración de NSGA-II se eligió con
validación cruzada de 5 pliegues agrupada por colegio y se eligió en 4 de 5; en el conjunto de prueba (32
listas, uso único) el hipervolumen fue 0,9962 frente a 0,9996 de la configuración anterior (diferencia
$-$0,0034, IC [$-$0,0054, $-$0,0017]), con un tiempo medio que bajó de 2,04 s a 0,47 s. El umbral de
enumeración exacta de 4.096 se respalda en que el exacto deja de ganar alrededor de $S \approx 4.932$
(mediana). En la sensibilidad de los pesos, el kit recomendado se mantiene en 0,965, 0,986 y 0,946 de los
sorteos para $\sigma$ de 0,05, 0,10 y 0,20. Sobre $\tau$ no se cita precisión, cobertura ni AUC, porque no
hay una calibración válida: el valor por omisión 0,75 queda como limitación declarada.

\subsection{Pruebas de carga}

\textbf{No ejecutadas en la primera corrida} (ver la sección~\ref{sec:reejecucion} para la línea base de lectura). No hay resultados de Locust, k6 ni de otra herramienta de carga en \ruta{ensayos/} ni
en \ruta{docs/}. Los requisitos DP-02 (50 usuarios concurrentes) y DP-05 quedan sin verificar.

\subsection{Usabilidad (OE5)}

Fuera del alcance de este informe. La ronda formativa fue el 31 de octubre y la sumativa está programada para
el 2 de noviembre de 2026, así que su resultado está \textbf{pendiente}. El requisito UI-07 (interacciones
discretas y RTLX, según la enmienda R1) depende de esa ronda.

\section{Verificación legal por inspección}

Esta sección es una \textbf{inspección}, no una prueba automatizada: se leyeron las páginas servidas y el
código contra la Ley 1581 de 2012 (con el Decreto 1377 de 2013) y la Ley 1480 de 2011. Donde existe una prueba
automatizada que respalda el punto se nombra.

\begin{longtable}{L{4.2cm}L{6.6cm}L{3.8cm}}
\caption{Lista de verificación legal (1 de noviembre de 2026).}\\
\toprule
Punto & Lo observado & Resultado \\
\midrule\endhead
Páginas legales servidas & \ruta{/\allowbreak{}legal}, \ruta{/\allowbreak{}legal/\allowbreak{}privacidad}, \ruta{/\allowbreak{}legal/\allowbreak{}terminos} y
\ruta{/\allowbreak{}legal/\allowbreak{}retracto} responden 200 & cumple \\
Política de datos & cita la Ley 1581 y el Decreto 1377, explica la supresión y los plazos de respuesta (10
días hábiles para consultas, 15 para reclamos) & cumple \\
Retención de datos de investigación & la política no menciona los 12 meses tras la aprobación del trabajo de
grado ni los 30 días de las grabaciones & observación: ese plazo vive en el consentimiento de investigación,
no en la política de la app \\
Términos & citan la Ley 1480 y la Ley 1581 & cumple \\
Página de retracto & cita la Ley 1480, los 5 días hábiles y los 30 días para el reembolso & cumple \\
Plazos en el dominio & \ruta{ley1480.py}: \ruta{DIAS\_\allowbreak{}HABILES\_\allowbreak{}RETRACTO = 5}, \ruta{DIAS\_\allowbreak{}REEMBOLSO = 30};
día hábil es lunes a viernes sin festivos de Colombia, con la ley Emiliani y la Semana Santa calculadas; el
plazo corre desde la entrega & cumple; respaldo: \ruta{test\_\allowbreak{}dominio\_\allowbreak{}ley1480.py} (7) y
\ruta{test\_\allowbreak{}retracto.py} (14), pasan \\
Consentimiento en el registro & la casilla de \ruta{/\allowbreak{}empezar} bloquea el envío si no se marca; el servicio de
identidad registra el consentimiento \ruta{registro\_\allowbreak{}titular} solo si \ruta{acepta\_\allowbreak{}terminos} es verdadero &
observación: el servidor no rechaza un registro sin aceptación, solo omite el registro del consentimiento \\
Consentimiento por finalidad y por hijo & revocación por hijo y por finalidad & cumple; respaldo:
\ruta{test\_\allowbreak{}privacidad.py} (14), pasa \\
Supresión & \ruta{POST /\allowbreak{}api/\allowbreak{}cuenta/\allowbreak{}supresion} sin sesión responde 401 (ejecutado); exige una frase de
confirmación; la solicitud no se completa hasta que confirman todos los servicios & cumple; respaldo:
\ruta{test\_\allowbreak{}supresion\_\allowbreak{}confirmada.py} (5) y las pruebas de supresión de comercio, inteligencia y operación \\
Exportación de datos & 200 en la E2E & cumple \\
Nombre del menor & el registro guarda solo un apodo opcional y cifrado, nunca el nombre legal; el intérprete
NL2C depura correo, teléfono, enlace y nombre & cumple; respaldo:
\ruta{test\_\allowbreak{}sanitizar\_\allowbreak{}quita\_\allowbreak{}correo\_\allowbreak{}telefono\_\allowbreak{}enlace\_\allowbreak{}y\_\allowbreak{}nombre} \\
\bottomrule
\end{longtable}

\section{Hallazgos}\label{sec:hallazgos}

Estado actualizado en la sección~\ref{sec:reejecucion}.

Los hallazgos se ordenan por severidad. Ninguno se corrigió en la primera corrida (1 de noviembre): la versión estaba congelada para la
ronda del OE5 y la regla de esas sesiones es registrar los defectos en lugar de parchear la app.

\begin{description}
  \item[H4 (media). Lectura de trabajos ajenos.] \ruta{GET /\allowbreak{}v1/\allowbreak{}trabajos/\allowbreak{}\{trabajo\_\allowbreak{}id\}}
  (\ruta{servicios/\allowbreak{}inteligencia/\allowbreak{}tukit\_\allowbreak{}inteligencia/\allowbreak{}routers/\allowbreak{}inteligencia.py:301}, función
  \ruta{ver\_\allowbreak{}trabajo}) no recibe la sesión ni compara al dueño. Cualquier
  persona con el identificador lee el kit completo, también sin sesión, y lo mismo pasa con las opciones de una
  línea. La web reenvía el token, pero el servicio no lo usa en esta ruta. Lo expuesto es la composición del kit, el colegio, el grado, el perfil y el personaje; no hay correo
  ni nombre. El identificador es un UUID versión 4 (122 bits aleatorios), así que no se puede enumerar, pero
  queda expuesto si el enlace se comparte. Las rutas de escritura del mismo trabajo sí comprueban al dueño. Esto
  contradice SG-02 en su lectura estricta.
  \item[H5 (media). Dependencias de la web.] \ruta{next} 16.3.4 con un aviso crítico y dos dependencias con
  avisos altos; subir a \ruta{next} 16.4.3 o posterior.
  \item[H3 (media). Posible token en el historial.] Dos coincidencias \ruta{github-pat} en archivos de
  \ruta{.opencode/\allowbreak{}handoff/\allowbreak{}history/} del commit \ruta{d0c19b0f}. Conviene revocar el token aunque haya vencido y
  decidir si se reescribe el historial del remoto.
  \item[H2 (baja). Humo desactualizado.] \ruta{arquitectura/\allowbreak{}humo.py} trae la clave de demostración en el código
  (línea 40) y pide la bandeja de operación sin token (línea 229). Debería leer la credencial del entorno y
  usar el token del operador que ya obtiene.
  \item[H1 (baja). Error de tipos.] \ruta{libs/\allowbreak{}ts/\allowbreak{}tukit-contrato/\allowbreak{}src/\allowbreak{}linea.ts:195}, TS2532: el resultado de
  \ruta{texto.split(/\allowbreak{}[:.;]/\allowbreak{}, 1)[0]} puede ser \ruta{undefined} según el compilador. En la práctica
  \ruta{split} siempre devuelve al menos un elemento y las 67 pruebas del paquete pasan, pero
  \ruta{tsc --noEmit} falla.
  \item[H6 (baja). Análisis estático.] 62 avisos de \ruta{ruff}, 16 archivos sin formato y 1 error de
  \ruta{eslint} en \ruta{AccionesPrivacidad.tsx:85}.
  \item[O1 (observación). Consentimiento solo en el cliente.] Ver la lista legal.
  \item[O2 (observación). Sin Service Worker.] La búsqueda de \ruta{serviceWorker} en \ruta{apps/\allowbreak{}web} no
  encontró registro de un Service Worker, lo que respalda marcar el offline-first como no verificado (enmienda
  R6).
\end{description}

\section{Trazabilidad hacia los requisitos}

El SRS v6 tiene 132 requisitos (67 RF, 49 RNF y 16 RML, según la hoja de hechos). La matriz cubre los que
tienen evidencia en este informe o una decisión de alcance, más los cuatro requisitos propuestos sin
identificador de la enmienda R5. El resto no tiene caso asociado aquí y no se debe leer como verificado. Los
números entre paréntesis son los casos que pasaron en cada archivo de prueba el 1 de noviembre.

\begin{longtable}{L{2.2cm}L{3.6cm}L{6.3cm}L{2.6cm}}
\caption{Matriz requisito, caso y resultado.}\\
\toprule
Requisito & Tema & Caso(s) & Resultado \\
\midrule\endhead
NR-01 (R4.a) & Estados de resolución & \ruta{test\_\allowbreak{}dominio\_\allowbreak{}pipeline} (12), \ruta{test\_\allowbreak{}trabajos} (52) & pasa \\
NR-04, ML-05 & Nada fuera del catálogo & \ruta{test\_\allowbreak{}validar\_\allowbreak{}linea} (10); personaje fuera del catálogo en
\ruta{test\_\allowbreak{}nl2c} & pasa \\
NR-07 & Reuso de lo validado & \ruta{test\_\allowbreak{}borrador\_\allowbreak{}canonico} (13), \ruta{test\_\allowbreak{}fastroute} (16); E2E
aplicar canónico & pasa \\
OP-01 (R4.b) & Kits por perfil & \ruta{test\_\allowbreak{}optimizador} (80); humo: 3 kits publicados & pasa \\
OP-02 & Parámetros documentados & artefactos de calibración & evidencia documental \\
OP-04, OP-07 (R4.c) & Pesos P21 y MAUT & \ruta{test\_\allowbreak{}optimizador}; análisis de sensibilidad & pasa \\
OP-05, ML-08 & Reproducibilidad & \ruta{test\_\allowbreak{}es\_\allowbreak{}reproducible\_\allowbreak{}bajo\_\allowbreak{}la\_\allowbreak{}misma\_\allowbreak{}semilla} & pasa \\
DP-02, DP-05 & Carga & no hay prueba de carga & no ejecutado \\
DP-03 & Tiempo del optimizador & la calibración mide tiempos medios, no el p95 del criterio & no verificado \\
AD-03, AD-06 & Revisión humana & \ruta{test\_\allowbreak{}admin\_\allowbreak{}sesion} (61), \ruta{test\_\allowbreak{}api\_\allowbreak{}inteligencia} (22); humo:
cola 401 sin token, aprobación y publicación & pasa \\
SG-01 & Claves con hash & \ruta{test\_\allowbreak{}dos\_\allowbreak{}hashes\_\allowbreak{}de\_\allowbreak{}la\_\allowbreak{}misma\_\allowbreak{}clave\_\allowbreak{}difieren\_\allowbreak{}por\_\allowbreak{}la\_\allowbreak{}sal} & pasa \\
SG-02 & Autorización con alcance & \ruta{test\_\allowbreak{}autorizacion} (19), \ruta{test\_\allowbreak{}mis\_\allowbreak{}casos} (7), pedido ajeno
en \ruta{test\_\allowbreak{}checkout}; IDOR ejecutado & \textbf{parcial}: H4 \\
CO-01, CO-02 & Panel de colegio & alcance reducido, no existe en el producto & no aplica \\
CO-03, CO-04 & Ámbito de colegio, matrícula & \ruta{test\_\allowbreak{}autorizacion} (19), \ruta{test\_\allowbreak{}api\_\allowbreak{}identidad} (70) & pasa \\
FU-01 & Orden de compra & \ruta{test\_\allowbreak{}checkout} (61); E2E ordenar & pasa \\
FU-03, UI-05 & Pago con Wompi & \ruta{test\_\allowbreak{}pasarela} (35); E2E con pago simulado; sandbox no ejecutado &
\textbf{parcial} \\
FU-04 (R5.1) & Retracto & \ruta{test\_\allowbreak{}retracto} (14), \ruta{test\_\allowbreak{}dominio\_\allowbreak{}ley1480} (7); E2E 409 y 201;
inspección & pasa \\
FU-08, DE-02 & Precio y entrega congelados & \ruta{test\_\allowbreak{}crear\_\allowbreak{}orden\_\allowbreak{}congela\_\allowbreak{}precio\_\allowbreak{}y\_\allowbreak{}reserva},
\ruta{test\_\allowbreak{}precio\_\allowbreak{}congelado} (2) & pasa \\
PV-01 & Consentimiento granular & \ruta{test\_\allowbreak{}privacidad} (14); inspección del registro & pasa, con O1 \\
PV-02 & Panel de privacidad & \ruta{test\_\allowbreak{}mis\_\allowbreak{}datos} (11); E2E exportar 200; supresión 401 sin sesión & pasa \\
PV-03, PE-06 & Datos del menor & depuración en \ruta{test\_\allowbreak{}nl2c}; inspección del registro & pasa \\
PV-04, PS-03 & Supresión completa & \ruta{test\_\allowbreak{}la\_\allowbreak{}solicitud\_\allowbreak{}no\_\allowbreak{}se\_\allowbreak{}completa\_\allowbreak{}hasta\_\allowbreak{}que\_\allowbreak{}confirmen\_\allowbreak{}todos};
pruebas de supresión de comercio (5), inteligencia (5) y operación (4) & pasa \\
MT-01 & Cobertura de 80~\% & sin \ruta{pytest-cov} & no ejecutado \\
MT-03 & Funciones de aptitud & \ruta{test\_\allowbreak{}fronteras} (47), \ruta{test\_\allowbreak{}fronteras\_\allowbreak{}se\_\allowbreak{}activan} (28) & pasa \\
UI-02, UI-06, PT-02 & Navegadores y dispositivos & sin pruebas en dispositivos & no ejecutado \\
UI-03 & Cambiar un ítem & \ruta{test\_\allowbreak{}descartar\_\allowbreak{}y\_\allowbreak{}opciones} (8); E2E carrusel y descarte & pasa \\
UI-07 (R1) & Carga cognitiva & ronda sumativa del OE5 & pendiente \\
R5.2 & Páginas legales & inspección & cumple, con observación de retención \\
R5.3 & NL2C & \ruta{test\_\allowbreak{}nl2c} (50), \ruta{test\_\allowbreak{}nl2c\_\allowbreak{}ensayo} (13); rondas 2 y 3 & pasa \\
R5.4 & Detalle de orden y logística admin & \ruta{test\_\allowbreak{}admin\_\allowbreak{}pedido} (6), \ruta{test\_\allowbreak{}logistica} (57);
pruebas TS de la consola admin & pasa \\
R6 & Offline-first & inspección: no hay Service Worker & no verificado \\
\bottomrule
\end{longtable}

\section{Lo que no se ejecutó}

\begin{itemize}
  \item Cobertura de código: \ruta{pytest-cov} no está instalado.
  \item Auditoría de dependencias de Python: \ruta{pip-audit} no está instalado.
  \item Pago real con el sandbox de Wompi: las llaves están pendientes.
  \item Pruebas de carga: no existen en el repositorio.
  \item Pruebas en navegadores y dispositivos reales (PT-02, UI-06): no se hicieron; Playwright no es estable en
        esta máquina.
  \item Prueba en vivo del intérprete NL2C dentro de la suite (\ruta{-m nl2c\_\allowbreak{}vivo}): reemplazada por las rondas
        del ensayo.
  \item Acceso cruzado sobre rutas de escritura y sobre el caso PQRS ajeno: fuera del alcance de una prueba de
        solo lectura; lo cubren pruebas automatizadas.
  \item Usabilidad sumativa del OE5: programada para el 2 de noviembre de 2026.
\end{itemize}

\section{Re-ejecución 2026-11-02}\label{sec:reejecucion}

Después de la primera corrida se corrigieron los hallazgos H1, H2, H4, H5 y H6 (parcial) antes de la segunda ronda
del OE5 y se repitieron las pruebas con el código ya confirmado en cuatro commits de la rama
\ruta{validacion/\allowbreak{}oe5} más uno de ajustes de revisión. Esta sección solo recoge resultados ejecutados el 2 de noviembre de 2026.
Se reiniciaron los servicios (\ruta{make reanudar}, con identidad antes que inteligencia) y se reconstruyó la web con
\ruta{next} 16.4.3.

\subsection{Commits}

\begin{itemize}
  \item \ruta{a201f4c}: higiene (\ruta{next} 16.4.3, tipos, lint, F811, clave fuera de \ruta{humo.py}, retención del estudio en la política de privacidad).
  \item \ruta{c147a6b}: cobertura con \ruta{pytest-cov}, ensayo de carga de lectura y barrido de supresiones atascadas.
  \item \ruta{0c9d1e1}: puerta de lectura del dueño en el kit (H4) y consentimiento obligatorio en el registro.
  \item \ruta{0059f6d}: \ruta{marcas\_\allowbreak{}excluidas} y \ruta{solo\_\allowbreak{}disponibles} como filtros duros del optimizador.
  \item \ruta{976deb1}: ajustes de revisión (credencial remota, carrusel, sondeo con 401, consentimiento en \ruta{/v1/cuenta}).
\end{itemize}

\subsection{Pruebas automatizadas}

\begin{table}[h]
\centering\small
\caption{Resultados de la re-ejecución (2 de noviembre de 2026).}
\begin{tabular}{L{5.2cm}L{4.6cm}rrr}
\toprule
Nivel & Comando & Pasan & Fallan & Saltadas \\
\midrule
Fronteras de arquitectura & \ruta{make fronteras} & 75 & 0 & 1 \\
Python completo & \ruta{make prueba} & 2.542 & 0 & 1 \\
Web del acudiente (TS) & \ruta{npm test} en \ruta{apps/\allowbreak{}web} & 86 & 0 & 0 \\
Contrato compartido (TS) & \ruta{npm test} en \ruta{libs/\allowbreak{}ts/\allowbreak{}tukit-contrato} & 67 & 0 & 0 \\
Consola admin (TS) & \ruta{node --test} (corrida del 2 de nov. por el track C) & 57 & 0 & 0 \\
Tipos TS & \ruta{tsc --noEmit} en los 3 paquetes & 3 de 3 & 0 & \\
\bottomrule
\end{tabular}
\end{table}

En la primera pasada de \ruta{make prueba} falló una prueba nueva (\ruta{test\_\allowbreak{}filtros\_\allowbreak{}duros.py}) porque importaba
\ruta{conftest} dentro de la función y tomaba el de operación. Se movió la importación al inicio del archivo y la
segunda pasada dio 0 fallas.

\subsection{Cobertura de código}

Medida con \ruta{pytest-cov} sobre los commits \ruta{a201f4c} a \ruta{0059f6d} (antes de los ajustes de \ruta{976deb1}); no hay umbral de falla.

\begin{table}[h]
\centering\small
\begin{tabular}{lr}
\toprule
Servicio & Cobertura \\
\midrule
identidad & 77 \% \\
catálogo & 54 \% \\
inteligencia & 88 \% \\
comercio & 83 \% \\
operación & 81 \% \\
gateway & 98 \% \\
\ruta{tukit\_comun} & 85 \% \\
\bottomrule
\end{tabular}
\end{table}

El 54 \% de catálogo viene de scripts de una sola corrida sin pruebas (\ruta{completar.py}, \ruta{vista.py},
\ruta{demanda\_\allowbreak{}atributos.py} y \ruta{representantes.py} con 0 \%).

\subsection{Pruebas de carga}

Ensayo de solo lectura (\ruta{ensayos/\allowbreak{}carga/\allowbreak{}carga.py}) directo a catálogo e inteligencia, 30 s por nivel, en un Apple M4
con el generador en la misma máquina. Error 0 \% en todos los niveles.

\begin{table}[h]
\centering\small
\begin{tabular}{rrrrr}
\toprule
Usuarios & Peticiones/s & p50 (ms) & p95 (ms) & p99 (ms) \\
\midrule
1 & 171 & 4,0 & 14,4 & 17,1 \\
5 & 889 & 5,0 & 9,6 & 11,6 \\
10 & 962 & 9,0 & 22,5 & 33,1 \\
20 & 891 & 18,9 & 48,7 & 71,3 \\
\bottomrule
\end{tabular}
\end{table}

Por el gateway, con 5 usuarios y 15 s, hubo 1.200 respuestas 2xx y 30.257 respuestas 429: es el limitador por IP
funcionando. Esto no verifica DP-02 completo: faltan los 50 concurrentes y los flujos de generar kits y checkout,
que escriben en la base. El rendimiento se satura entre 5 y 10 usuarios y el p95 sube de 9,6 a 48,7 ms entre 5 y 20.

\subsection{Seguridad}

\begin{itemize}
  \item \ruta{npm audit --omit=dev}: 0 vulnerabilidades (antes 3: 1 crítica y 2 altas).
  \item IDOR (\ruta{idor.sh}, contra \ruta{127.0.0.1:3000}), cuenta atacante: trabajo ajeno 404, opciones de línea ajena 404, pedido ajeno 404, trabajo ajeno sin cookie 401, pedido ajeno sin cookie 401, pedido inventado 404; casos propios 200 con 0 elementos y exportación propia sin identificadores ajenos.
  \item Dueño del trabajo con su cookie (por el túnel): API del trabajo 200, opciones 200, página del kit 200; la misma API sin cookie da 401.
  \item Enlace rápido sin clave (\ruta{/empezar}, ruta \ruta{aplicar-canonico} sin cuenta): API 200 y página \ruta{/kits/<id>} 200 sin sesión.
\end{itemize}

\subsection{Humo y extremo a extremo}

\begin{itemize}
  \item \ruta{make humo} (clave desde el llavero, sin imprimirla): 20 comprobaciones, todas con datos reales; la bandeja de operación se pide con token de operador.
  \item E2E del acudiente por el túnel (\ruta{e2e.sh}, \ruta{e2e2.sh}, \ruta{e2e3.sh}): registro con cookie, kit publicado con 3 perfiles y 72 líneas sin precios en cero, opciones (42), propuesta del asistente, orden, descarte (PUT 200; el POST de \ruta{e2e2.sh} da 405, es esperado), caso 201, exportación 200 y las seis páginas 200. Dos pasos no avanzan porque el pago no se completa en sandbox: \ruta{sincronizar} da 404 \ruta{transaccion\_\allowbreak{}no\_\allowbreak{}existe} y \ruta{retracto} da 409 \ruta{retracto\_\allowbreak{}sin\_\allowbreak{}pagar}; el pedido queda en \ruta{pendiente\_\allowbreak{}pago}.
\end{itemize}

\subsection{Estado de los hallazgos}

\begin{table}[h]
\centering\small
\begin{tabular}{llL{8.2cm}}
\toprule
Hallazgo & Estado & Evidencia \\
\midrule
H1 & resuelto & \ruta{a201f4c}; \ruta{tsc --noEmit} limpio en los 3 paquetes \\
H2 & resuelto & \ruta{a201f4c}; \ruta{make humo} lee la clave del llavero y pasa con 20 comprobaciones \\
H3 & abierto & token en el historial: requiere revocarlo y decidir sobre el remoto \\
H4 & resuelto & \ruta{0c9d1e1} y \ruta{976deb1}; IDOR: 404 o 401 ajeno, 200 el dueño \\
H5 & resuelto & \ruta{a201f4c}; \ruta{next} 16.4.3, \ruta{npm audit} en 0 \\
H6 & parcial & eslint de admin sin avisos y web con 0 errores y 1 aviso; \ruta{ruff} baja de 62 a 51 avisos y quedan 14 archivos sin formato \\
\bottomrule
\end{tabular}
\end{table}

Queda abierto además: \ruta{next} 16.4.4 no se pudo instalar por la política de antigüedad mínima de npm; el barrido de supresiones
se probó solo sobre SQLite en memoria y no se lanzó contra la base viva ni se añadió a \ruta{make arriba}; el reenvío del Bearer en
\ruta{cliente-http.ts} se cubre con \ruta{tsc} y revisión, sin prueba de ejecución.

\section{Evidencia}

La carpeta \ruta{QC/\allowbreak{}evidencia/} guarda los registros de cada corrida: \ruta{pytest\_\allowbreak{}prueba.log},
\ruta{pytest\_\allowbreak{}por\_\allowbreak{}servicio.txt}, \ruta{pytest\_\allowbreak{}por\_\allowbreak{}archivo.tsv}, \ruta{pytest\_\allowbreak{}colas\_\allowbreak{}redis.log}, los
resúmenes de TypeScript y de \ruta{tsc}, \ruta{lint.log}, \ruta{eslint.log}, los tres registros de humo, el
guion y la salida de IDOR, \ruta{legal\_\allowbreak{}paginas.txt}, los resúmenes de \ruta{npm audit} y los conteos de
\ruta{gitleaks}, estos sin ningún valor de secreto.

\end{document}
